\section{迁徙把王朝的边界带进家门}
\label{sec:synthesis-migration-households}

永嘉五年，洛阳的宫门失守时，逃离的人并没有沿着一条由史书替他们画好的南渡线移动。有人先到许昌、襄阳或江陵，等待渡口、亲属、船只和能够携带的粮食；有人留在原郡县，改向新来的政权纳粮服役；有人在坞壁、山谷和城郭之间反复迁居。朝廷后来需要把这些凌乱的去向写成郡、县、户、籍和官职，才可能征调、赈济、授官或断狱。因而迁徙并非西晋覆亡之后才附在正统叙事旁的一页悲剧，而是东晋、刘宋、北魏、北齐、北周和隋不断遇到的行政难题：一个人已经到了哪里，他仍然算谁的户，他的劳役、田土、婚姻和旧日依附应当由谁承认？

《晋书》将洛阳陷落、怀帝被俘与政权瓦解排入清晰的年序，使我们知道都城秩序何时被击穿；但它不能列出每一户人的路线。地方志、墓志和侨州郡名目把若干人钉在新的地点，反倒显示移动不是一次性完成的。墓石上的本贯常常保留北方郡望，葬地和仕宦却在江南；这种并置不是简单的“念旧”，而是迁居者在新秩序中争取资格的方式。它能说明一户有能力留下文字的家庭怎样安排身份，不能把无名船夫、佃客、被掠者和途中死亡者一并变成士族故事。

\subsection{建康不是南渡的终点}

司马睿在建康立足，首先面对的是一座并未为北方朝廷预先腾空的城市。长江下游原有孙吴以来的州郡、豪强、仓廪、渡口和地方关系；北来官员所携带的皇统名义需要落到这些既存网络上。新朝可以任命侨寓的北方人士，也须依赖江东地方领袖维持征收、治安与船运。把东晋写成“衣冠南渡”固然抓住了都城文化和部分精英的移动，却会抹掉一项更艰难的工作：如何让两套籍贯、声望和资源要求在同一片水网中不至于立刻相撞。

江面在这里不是天然的屏障。它既能阻挡北方军队，也要求朝廷掌握船只、港埠、堤岸与熟悉潮汐的人。荆州连接上游，京口与广陵望向淮北，豫章、会稽和闽岭方向则关联粮食、山货与海路。一个从北方来的人若只抵达建康城外，尚未获得寄居处、食物和能被官府承认的身份；一个本地人若因征发、战争或水患离开乡里，也未必因此进入“北来者”这个类别。迁徙的政治意义恰在于这些不同的路线在州郡档案和市场中相遇。

\begin{turningpoint}[永嘉之后：皇统南移，登记问题没有随之南移完成]
元帝在江左建立朝廷，给许多北来者提供了可以宣称的政治中心，却没有自动提供一张完整户口册。侨置州郡、寄籍、授官和赈济都在尝试把失散的人编入新的责任链。它们解决的是“谁可以被看见、由谁负担”的当下问题，代价是同一空间里可能并存旧县、侨县与实际居住地，后来的征税、服役和诉讼必须穿过这些重叠名称。
\end{turningpoint}

这也是为什么东晋政局的每一次北伐都会同时触动江南社会。祖逖在河南的经营、桓温和刘裕的北上，不只是地图上向北推出一段颜色；军队需要从江淮与荆州取得粮运，前线收复的城镇又会吸引流民、旧吏和地方武装重新选择归属。北伐失败时，先前依附的家户可能再次南移，或被留在新占领区之外。刘裕军进入洛阳和关中时，朝廷短暂拥有重新连接旧都与南方的想象；其军队撤离后，水陆补给和地方接收的断裂也同样真实。胜败因此改变的不是抽象的人口箭头，而是具体家庭能否把住处、职业和身份安放在一个可持续的政权之下。

\subsection{侨州郡：名称可以保存来处，不能替代住处}

侨置州郡最容易被读成一份精确的流民统计。名称沿用旧郡县，似乎每一个“侨”字都对应一批由原籍整齐迁来的居民；其实行政名称的保留首先服务于官职、名望、赋役归属和政治承认。它使北方旧地在江南官制里继续可见，也让朝廷能够对迁来的人说：你并非无籍之人。但同一个侨置名目所覆盖的人口、土地和实际治所可以变化；有人被编入其中，有人只是暂住，有人继续按原有地方网络生活。若把名目直接换算为人数，便会把国家想要整理的秩序误当成已经实现的秩序。

东晋与南朝的“土断”更能说明这种差别。把寄居者依其实际住地编入当地户籍，表面上是清理名籍，实际触到的是谁承担赋役、谁失去旧有保护、地方官能否掌握人丁的问题。一次土断并不是简单的行政进步：对朝廷而言，它可能增加可征资源；对地方强宗和侨寓精英而言，它可能改变依附者、佃客与婚姻网络被承认的方式；对普通家户而言，则意味着在新的里、县和差役序列中重新出现。史书记录改革者的意图和成败，较难留下被重新登记的人怎样理解这件事。我们应当承认这个沉默，而不把“编户”写成所有人欣然接受的稳定状态。

南朝都城附近的墓志有时可使这种抽象过程变得可见。一方墓志所列的本贯、官职、婚配和葬地，能让读者看见某个家族把北方郡望同江南仕宦联结起来；它也提醒我们，墓志是后人选择保留的名誉文本。文中没有列出的乳母、佣作、依附者和女方亲属，并不因此不存在。以墓志读家庭，最有价值的并非替所有人填出一张家谱，而是让“迁徙”恢复为不同成员承担不同风险的过程：能出仕者要取得名分，留守者要保有田产，婚配要使两地关系可被信任，葬地则把一段新生活固定在具体土壤中。

\begin{causalchain}[侨置与土断的制度得失]
侨置使朝廷能在北方失陷后继续承认旧籍、安置官员并维持皇统的地理记忆；它也让实际居住、旧日身份和赋役责任可以彼此错开。土断试图把这种错开重新压回地方登记，以便征收和治理，却可能削弱既有保护关系并引起抵触。二者不是一进一退的抽象制度史，而是在战争、粮运和地方竞争中反复调整的工具。
\end{causalchain}

\subsection{北方的移动不是一条“胡人南下”}

北方同样有大规模移动，但它们不能被一个族名概括。西晋末年，关中、并州、河北和河东的城镇、坞壁与牧地承受不同压力；前赵、后赵、前燕、前秦等政权分别倚重军队、部众、旧郡县吏、商路和城郭。史书中“降”“叛”“徙”“部”的词语往往写在军事行动之后，容易使读者以为一纸命令便带走了完整的人群。更实际的情形是，战争胜者必须决定俘获者、降者、旧吏与工匠怎样安置，失败者则要选择投奔哪座城、哪条河谷或哪一位地方首领。这个选择经常随着粮食、亲属与城门是否还开着而改变。

冉闵在邺城称帝和后赵解体，是对这种危险的极端提醒。载记和编年材料能够确认内战暴力与大规模逃散的发生，却无法把被害者、逃亡者和各地参与者精确分为后世想要的稳定族群。把“羯”“胡”“汉”直接当成固定民族名单，会遮住当时军政动员与排斥语言怎样制造新的危险。对家庭而言，最迫切的未必是接受某个宏大身份，而是避开城市清洗、保存亲属、找到食物和进入另一套能给保护的关系。北方人口的去向因此既是政治史，也是城防、道路与仓储史。

前秦一度整合关中、河北与河西，又在淝水后迅速崩散，正好显示“统一”与“可留住人”不是同一件事。前秦的任命、征发和道路控制能把多地暂时接在一起，但当军队失利、中心无法继续供给时，原先被压住的地方竞争、继承冲突和生计顾虑重新显形。吕光据凉州、各凉政权相继出现，并不只是雄主失败后的尾声；河西城镇、绿洲和商道提供了可重新组织的资源，因而使人群、军队与书写在北方碎裂后仍能寻找新的节点。

\subsection{从平城到洛阳：迁都也是人口与技能的调配}

拓跋珪迁都平城，表明一个以代北为重心的政权必须把军政中枢安放在可经营的城郭里。平城附近的道路、军镇、牧地和新吸纳的河北资源，使北魏可以同时面对内部整合与北方边患。迁都不应被写成从“游牧”到“农业”的单向跃迁：城内行政需要文书与官员，城外军队仍需马匹、草场和边地补给；旧有部众、降附者和州郡居民也未因此进入同一生活方式。

孝文帝迁都洛阳把这种调配推得更远。洛阳具有旧都的礼制与交通意义，也更接近中原郡县、运河和可征收的农业腹地；但离开平城意味着北方军人、边镇和原有政治网络要重新被安置和解释。朝廷倡导改姓、服饰与礼仪，确实塑造了公开可见的规范，却不能量出每一户的语言和自我认同。迁都的物质面同样重要：宫城、坊市、寺院、道路、桥梁和仓储需要劳力、木石与运输，迁入者则要在新城找到住处、职位和信任链。正是在这一层面，洛阳既是王朝的宣言，也是无数资源被集中和再分配的工地。

六镇问题暴露了这项调配的反面。边镇军人及其家属不只是皇帝在地图边缘摆放的棋子；他们的给养、婚配、升迁和迁移机会关系到是否愿意继续驻守。迁都后中枢的文化与仕宦重心改变，边镇的不满不能归结为单一的“鲜卑反汉化”或“汉人反鲜卑”。起事的时间、人物与后续军阀化可由史书追出，但不同镇户的利益和感受不能被一条原因替代。战争把这些家庭推入河北、关中与河东，日后东西魏、北齐和北周的军事与地方网络，都要处理由此生成的流动人口。

\subsection{江陵、邺与建康：失都以后谁带走了什么}

侯景之乱和江陵陷落使南方读者再次看见迁徙的暴力面。城市被围困时，离开并不只是一张船票：粮食先被军队、官府和富户争夺，城门、江岸与道路决定哪些人能够走，留下的人则可能被征发、掠夺或迫使改投新主。后出的文人记忆尤其擅长保存失都者的悲痛，但这种记忆不能替全部市民作证。我们仍可从它们提出更具体的问题：书籍、官署文书、工匠技术和亲族关系怎样离开旧城；新政权为什么需要接收其中的一部分；又有哪些人的损失没有进入书写。

北方的邺城在东魏、北齐时期积聚军队、官僚、工匠与佛教供养网络，北周灭齐后同样面临接收与再编的难题。城市并不会因改朝换代而立即失去原有道路、市场和熟练劳力；新政权常常必须利用这些遗产，同时防范旧军队、旧贵族或地方势力。隋的再统一也应在这个意义上理解：五八九年陈亡使最高层的敌对朝廷消失，但江南家户、寺院、州郡与水路网络并不会在同一天获得相同的新身份。再统一的速度，要用接收、登记、供给和裁决能否在地方持续来衡量。

\begin{debate}[南渡人数能算出来吗]
传世史书、侨置记录和墓志都表明北方人口持续进入江南，但它们不是同一种计量材料。名门墓志偏向保存能立碑的人，侨置名目服务行政承认，正史叙事则偏爱危机与名士。现代研究可以比较这些线索、估计迁徙对江南社会的影响，却不宜从任何一类材料推出一个无争议的总数，或以一条总数解释所有地区和阶层的经历。
\end{debate}

迁徙的长时段后果因而不是“北人带来文化、南人接受文化”的单线故事。北方旧籍在江南继续提供名望，江南土地、水路和地方关系也改变了这些人的生活；平城、洛阳、邺和长安则吸纳来自不同方向的军人、官员、工匠、僧侣与商旅。家庭把这些变动记在婚姻、教育、墓地、债务和职业里，国家则试图把它们改写为可以计算的户、役、兵和田。两种记录永远不能完全重合；正是它们之间的缝隙，使两晋南北朝的政治更替不止发生在宫门之内。
